% R. Brown's translations of two letters to L. Breen, with his note and comments (cote 134-2): archive cote 134-2, site page 124 (PDF page 125).
% translation of the letter of 17/2/1975. Draft compared: none.
% AI-transcribed by Claude Opus 5.5 (claude-opus-5-5, adaptive thinking, effort high), Message Batches API, batch pilotnext (msgbatch_01A9sdu8SYq23bEtqX9tUncQ), 2026-11-13; not yet checked by eye.
\scanpage{134-2/124}
\mspage{- 3 -}

(i = 1, 2) and an equivalence between the restrictions of $F_1$ and $F_2$
to $T_1 \cap T_2$ (while the analogous statement in terms of 0-champs, i.e. for
coverings, is evidently correct).

The statement (B) above makes it clear how to give explicitly the
cohomology of an n-groupoid $C_n$ . If $C_n = \Pi_n(X)$ , and if $F$ is a locally
constant (n-1)-champ on $X$ , and $e^{X}_{n-1}$ is the ``final'' (n-1)-champ, one has
an (n-1)-equivalence of (n-1)-categories
\[ \Gamma_{\unsure{x|\times}}(F) = F(X) \simeq \mathit{Hom}(e^{X}_{n-1}, F) \]
which shows that the functor $\Gamma_X$ ``integration on $X$'' for locally constant
(n-1)-champs, which includes the (non-commutative) locally constant cohomology
of $X$ in dim $\leq$n-1, can be interpreted in terms of ``local (n-1)-systems''
on the fundamental groupoid as an
$\mathit{Hom}(e^{(C_n)}_{n-1} , F)$ where now $F$ is interpreted as an n-functor
\[ C_n \xrightarrow{\ F\ } ((n-1)\text{-Cat}) \]
and $e^{(C_n)}_{n-1}$ is the constant n-functor on $\unsure{\underline{C}_n|C_n}$ , with value the final (n-1)-
category.

To interpret this in cohomology notation, it is necessary for me to add,
as ``apology'' to the preceding letter, the explicit interpretation of the
non-commutative cohomology on a topos $X$ , in terms of integration of n-
champs on $X$ . If $F$ is a strict Picard n-champ on $X$ , then it is defined
by a complex $L^{\circ}$ of cochains on $X$
\[ 0 \to L^0 \to L^1 \to L^2 \to \ldots \to L^n \to 0 , \]
concentrated in degrees $0 \leq i \leq n$ (defined uniquely up to isomorphism in
the derived category of Ab(X)). That said, the $H^i(X, L^{\circ})$ (hypercohomology)
for $0 \leq i \leq n$ can be interpreted as $H^i(X, L^{\circ}) = \pi_{n-i}\Gamma_X(F)$ . If one is
interested in all the $H^i$ (not just for $i \leq n$) one must, for all $N \geq n$ ,
regard $L^{\circ}$ as a complex concentrated in degrees $0 \leq i \leq N$ (by prolongation
of $L^{\circ}$ by 0 to the right). The corresponding strict Picard n-champ is no
longer $F$ but $\mathbb{C}^{N-n}F$ , where $\mathbb{C}$ is the ``classifying space'' functor, interpreted
on strict Picard n-categories as the operation consisting of ``translating''
the i-objects to (i+1)-objects, and adjoining a unique 0-object; this extends
one hopes, in ``an obvious way'', to n-champs, so as to commute with the operation
of taking the inverse image of an n-champs. One has then for $i \leq N$
\[ H^i(X, L^{\circ}) = \pi_{N-i}\Gamma_X(\mathbb{C}^{N-n}F) \qquad i \leq N . \]
Given this, it is necessary to put, for all strict Picard n-champs $F$ on $X$ ,
\[ H^i(X, F) = \pi_{N-i}\Gamma_X(\mathbb{C}^{N-n}F) \quad \text{if} \quad N \leq i, n \]
which does not depend on the choice of integer $N \geq \sup(i, n)$ .

[\uline{N.B.} One has a canonical morphism of (n-1)-groupoids
\[ \mathbb{C}(\Gamma_X F) \longrightarrow \Gamma_X(\mathbb{C}F) , \]
as the obvious constructions in terms of cochains show, and one sees in the

\mspage{\hand{125}}
